\documentclass[11pt]{article}
\usepackage[margin=1in]{geometry}
\usepackage{amsmath,amssymb,amsthm,mathtools,mathrsfs}
\usepackage[colorlinks=true,allcolors=blue]{hyperref}
\newtheorem{theorem}{Theorem}
\newtheorem{lemma}[theorem]{Lemma}
\newtheorem{proposition}[theorem]{Proposition}
\newtheorem{corollary}[theorem]{Corollary}
\newcommand{\B}{\mathscr B}
\newcommand{\Sch}{\mathcal S}
\newcommand{\HH}{\mathcal H}
\newcommand{\Z}{\mathbb Z}
\newcommand{\R}{\mathbb R}
\newcommand{\C}{\mathbb C}
\newcommand{\N}{\mathbb N}
\newcommand{\id}{\operatorname{id}}
\newcommand{\cl}{\operatorname{cl}}
\newcommand{\supp}{\operatorname{supp}}
\newcommand{\F}{\mathcal F}
\title{The closed range of the original full theta source\\
with every weight, derivative, and factor retained}
\author{}
\date{20 September 2026}
\begin{document}
\maketitle

\section{The fixed objects and the source theorem}

Every space and map in this note is over $\C$. The original source is
\[
 V=\left\{\phi\in\Sch(\R):\phi(-x)=\phi(x),\quad
       \phi(0)=0,\quad \int_\R\phi(x)\,dx=0\right\},
\]
with the subspace topology from the ordinary Schwartz space. Put
\[
 D=-x\frac{d}{dx},\qquad
 p_{b,j}(F)=\sup_{x>0}x^b|D^jF(x)|\quad(b\in\Z,\ j\in\N_0),
\]
\[
 \B=\{F\in C^\infty(\R_{>0}):p_{b,j}(F)<\infty
                     \text{ for every }b,j\}.
\]
The topology on $\B$ is exactly that generated by all $p_{b,j}$. The
operators and Fourier convention are
\begin{equation}\label{eq:original}
 \Theta\phi(x)=2\sum_{n=1}^{\infty}\phi(nx),\qquad
 JF(x)=x^{-1}F(x^{-1}),\qquad
 \widehat\phi(\xi)=\int_\R\phi(x)e^{-2\pi i x\xi}\,dx.
\end{equation}
No conjugation by $x^{1/2}$, renormalization of $\Theta$, change of the
coordinate $x$, or substitution of a Hilbert norm is made here.

The literature source is Ralf Meyer, \emph{A spectral interpretation for
the zeros of the Riemann zeta function}, arXiv:math/0412277v3,
\href{https://arxiv.org/abs/math/0412277v3}{original versioned source}.
Its Section 2 defines
\[
 \HH_- =\bigcap_{s\in\R}\Sch(\R_+^*)_s,\qquad
 \Sch(\R_+^*)_s=
 \{F:x\mapsto x^sF(x)\text{ is Schwartz in }\log x\}.
\]
Its Section 3 defines $Z\phi(x)=\sum_{n\geq1}\phi(nx)$ and
$\HH_\cap=\{\phi\in\Sch(\R):\phi\text{ even},\,
\phi(0)=\F\phi(0)=0\}$. Theorem 3.3, with source label
\texttt{the:Zeta\_estimate}, proves closed range and the
topological-isomorphism-onto-range property, including the restriction
$Z:\HH_\cap\to\HH_-$. Its proof uses Poisson summation and
Proposition 3.2 (label \texttt{pro:Z\_invertible}), together with
Proposition 2.1 (label \texttt{pro:Sch\_estimates}).

Meyer's Fourier kernel has the opposite sign. For an even $\phi$ the
change of variable $x\mapsto-x$ in the defining integral proves that
the two transforms are equal, as functions, with no scalar factor.
Consequently $V=\HH_\cap$ with identical source topology. His $Z$ and
the original map satisfy $\Theta=2Z$. The next sections prove the target
topology identification and independently reconstruct the closed-range
conclusion with explicit inverse estimates.

\section{Identity equivalence of the complete seminorm systems}

For $s\in\R$ and $k,j\in\N_0$ define
\begin{equation}\label{eq:q}
 q_{s,k,j}(F)=\sup_{x>0}(1+|\log x|)^k
 \left|\left(x\frac d{dx}\right)^j(x^sF(x))\right|.
\end{equation}
These are the ordinary Schwartz seminorms in the variable $\log x$
after multiplication by $x^s$, expressed on the original $x$-axis. Their
projective family generates precisely Meyer's topology on $\HH_-$.

\begin{theorem}\label{thm:topology}
The identity on positive-axis functions is a topological linear
isomorphism $\B=\HH_-$. More explicitly, set
\[
 n_-(s)=\lfloor s\rfloor-1,\qquad n_+(s)=\lceil s\rceil+1,
 \qquad C_k=\sup_{t\geq0}(1+t)^ke^{-t}.
\]
Here $C_0=1$ and $C_k=k^ke^{1-k}$ for integer $k\geq1$.
For every real $s$ and every $k,j\geq0$,
\begin{equation}\label{eq:q-from-p}
 q_{s,k,j}(F)\leq C_k\sum_{\ell=0}^j\binom j\ell
 |s|^{j-\ell}\bigl(p_{n_-(s),\ell}(F)+p_{n_+(s),\ell}(F)\bigr).
\end{equation}
Conversely, for every integer $b$ and $j\geq0$,
\begin{equation}\label{eq:p-from-q}
 p_{b,j}(F)\leq\sum_{\ell=0}^j\binom j\ell
 |b|^{j-\ell}q_{b,0,\ell}(F).
\end{equation}
In these formulas a zeroth power, including $0^0$, means $1$.
\end{theorem}

\begin{proof}
Write $L=x\,d/dx=-D$. The Leibniz rule, iterated $j$ times, gives the
exact operator identity
\[
 L^j(x^sF)=x^s(s-D)^jF
   =x^s\sum_{\ell=0}^j\binom j\ell s^{j-\ell}(-1)^\ell D^\ell F.
\]
For $x\geq1$ put $t=\log x$. Since $n_+(s)-s\geq1$,
\[
 (1+|\log x|)^kx^s\leq C_kx^{n_+(s)}.
\]
For $0<x\leq1$ put $t=-\log x$. Since $s-n_-(s)\geq1$,
\[
 (1+|\log x|)^kx^s\leq C_kx^{n_-(s)}.
\]
Substitution in the displayed binomial expansion proves
\eqref{eq:q-from-p}. The constant follows by maximizing
$k\log(1+t)-t$ on $t\geq0$.

The conjugate identity $(b-L)(x^bF)=x^bDF$ yields
\[
 x^bD^jF=(b-L)^j(x^bF).
\]
Its binomial expansion gives \eqref{eq:p-from-q}. Hence finiteness of
the full $p$-family and of the full $q$-family is equivalent, and each
seminorm of either family is bounded by a finite linear combination
from the other. This proves equality of the sets and of their locally
convex topologies, through the identity map.
\end{proof}

For completeness, $\B$ is a Fr\'echet space directly in the $p$-topology.
The seminorm family is countable and separates points. If $(F_n)$ is
Cauchy in every $p_{b,j}$, it is Cauchy in $C^\infty(K)$ for each compact
$K\subset(0,\infty)$, because
\begin{equation}\label{eq:R}
 x^m\partial_x^m=R_m(D),\qquad
 R_m(t)=(-1)^m\prod_{r=0}^{m-1}(t+r),\qquad R_0(t)=1.
\end{equation}
The identity follows by induction, using
$x^{m+1}\partial^{m+1}=(x\partial-m)x^m\partial^m$.
Local limits therefore give a single smooth $F$. For a fixed $b,j$,
the $p_{b,j}$-Cauchy bound on $F_n-F_m$ passes pointwise to $m\to\infty$
and then to the supremum; thus $p_{b,j}(F_n-F)\to0$ and $F\in\B$.
The same local-limit argument proves completeness of $\Sch(\R)$.
Parity, evaluation at zero, and integration are continuous Schwartz
functionals/conditions (for integration use the seminorm with weight
$(1+|x|)^2$). Therefore $V$ is a closed, complete Fr\'echet subspace.

\section{Continuity, Poisson symmetry, and exact M\"obius inversion}

Write
\[
 T_{N,j}(\psi)=\sup_{t\geq1}t^N|D^j\psi(t)|
 \qquad(\psi\in\Sch(\R)).
\]
Each $T_{N,j}$ is a continuous Schwartz seminorm: expanding
$D^j=(-x\partial_x)^j$ gives a finite sum of polynomial coefficients
times ordinary derivatives. For $N>1$, termwise differentiation of
the theta series on every compact subset of $(0,\infty)$ is justified
by the convergent majorant $C n^{-N}$. On $x\geq1$ it gives
\begin{equation}\label{eq:theta-tail}
 |D^j\Theta\psi(x)|\leq2\zeta(N)x^{-N}T_{N,j}(\psi).
\end{equation}

The ordinary Fourier transform maps $\Sch(\R)$ continuously to itself.
Indeed, its differentiated, polynomially weighted expressions are
Fourier transforms of finite sums of derivatives of polynomially
weighted $\psi$, multiplied by the explicit powers of $2\pi i$ from
the kernel in \eqref{eq:original}; their suprema are bounded by the
corresponding $L^1$ norms. Those norms are bounded by finitely many
Schwartz seminorms, using the integrable weight $(1+|x|)^{-2}$.
Fourier inversion gives $\widehat{\widehat\psi}(x)=\psi(-x)$.
It follows that the Fourier transform is an involution of $V$:
evenness is preserved, $\widehat\phi(0)=\int\phi=0$, and
$\int\widehat\phi=\phi(0)=0$.

Poisson summation with the exact kernel \eqref{eq:original} says
\[
 \sum_{n\in\Z}\psi(nx)
    =x^{-1}\sum_{n\in\Z}\widehat\psi(n/x).
\]
It follows, for example, by periodizing $t\mapsto\psi(xt)$:
its $k$th Fourier-series coefficient is $x^{-1}\widehat\psi(k/x)$,
and the series and all needed derivatives converge absolutely by
Schwartz decay; evaluation at $t=0$ gives the formula.
For even $\psi$ this is exactly
\begin{equation}\label{eq:poisson}
 \psi(0)+\Theta\psi(x)
       =x^{-1}\widehat\psi(0)+J\Theta\widehat\psi(x).
\end{equation}
Both endpoint terms vanish for $\phi\in V$, so
\begin{equation}\label{eq:symmetry}
 \Theta\phi=J\Theta\widehat\phi,
 \qquad J\Theta\phi=\Theta\widehat\phi.
\end{equation}
The second identity uses $J^2=\id$.

Direct differentiation on the original axis gives
\begin{equation}\label{eq:J}
 D^jJ=J(1-D)^j,\qquad
 p_{b,j}(JF)\leq\sum_{r=0}^j\binom jr p_{1-b,r}(F).
\end{equation}
For $j=1$ the first identity follows from
$-x\partial_x(x^{-1}F(x^{-1}))
=x^{-1}F(x^{-1})+x^{-2}F'(x^{-1})$;
induction proves the general case. In the seminorm bound substitute
$y=x^{-1}$, so $x^b x^{-1}=y^{1-b}$, and expand $(1-D)^j$.
In particular $J$ is a continuous involution of $\B$.

Choose the integers $N=\max(2,b)$ and $N'=\max(2,1-b)$.
Combining \eqref{eq:theta-tail} on $x\geq1$ with
\eqref{eq:symmetry} and \eqref{eq:J} on $0<x\leq1$ yields
\begin{equation}\label{eq:theta-continuity}
 p_{b,j}(\Theta\phi)
 \leq2\zeta(N)T_{N,j}(\phi)
  +2\zeta(N')\sum_{r=0}^j\binom jrT_{N',r}(\widehat\phi).
\end{equation}
Thus $\Theta:V\to\B$ is well defined and continuous in the precise
original topologies.

For a positive-axis smooth function define
$p^+_{u,j}(\psi)=\sup_{x>0}x^u|D^j\psi(x)|$ whenever finite;
the superscript distinguishes a source estimate from the target
seminorm notation. Let $\mu(n)$ be the M\"obius function and put
\[
 (MG)(x)=\sum_{n=1}^{\infty}\mu(n)G(nx).
\]
Whenever $p^+_{u,r}(G)<\infty$ for $u>1$ and $0\leq r\leq j$, its
series may be differentiated $j$ times locally uniformly and
\begin{equation}\label{eq:mobius-bound}
 p^+_{u,j}(MG)\leq\zeta(u)p^+_{u,j}(G).
\end{equation}
Here $|\mu(n)|\leq1$ and $D^j[G(nx)]=(D^jG)(nx)$ give the estimate
term by term. In particular this applies to every $G\in\B$ and
integer $u\geq2$.

For $\phi\in V$ and fixed $x>0$, the double series
\[
 \sum_{n,m\geq1}|\mu(n)\phi(mnx)|
 \leq x^{-u}p^+_{u,0}(\phi)\zeta(u)^2
 \qquad(u>1)
\]
is finite. Grouping by $k=mn$ and using
$\sum_{n\mid k}\mu(n)=1$ for $k=1$ and $0$ otherwise gives
\begin{equation}\label{eq:mobius-inversion}
 M\Theta\phi=2\phi\quad\text{on }(0,\infty).
\end{equation}
The divisor identity follows by taking the product
$\prod_{p\mid k}(1-1)$ when $k>1$; for $k=1$ the empty product is $1$.
Thus $\Theta$ is injective. Writing $F=\Theta\phi$, the exact inverse
formulas and their estimates are
\begin{align}
 \phi(x)&=\tfrac12\sum_{n\geq1}\mu(n)F(nx),&
 p^+_{u,j}(\phi)&\leq\tfrac12\zeta(u)p_{u,j}(F),
       \label{eq:inverse-phi}\\
 \widehat\phi(x)&=\tfrac12\sum_{n\geq1}\mu(n)(JF)(nx),&
 p^+_{u,j}(\widehat\phi)&\leq\tfrac12\zeta(u)
                    \sum_{\ell=0}^j\binom j\ell p_{1-u,\ell}(F),
       \label{eq:inverse-hat}
\end{align}
for $x>0$, integers $u\geq2$, and $j\geq0$. The second line uses
\eqref{eq:symmetry}. These are inverse formulas on the image; they
do not assert that $MG$ extends to an element of $V$ for an arbitrary
$G\in\B$.

\section{Explicit recovery of every source Schwartz seminorm}

The following estimate supplies the part of the inverse-continuity
argument where both the function and its Fourier transform are needed.
All $L^2$ norms in this section use ordinary Lebesgue measure $dx$ or
$d\xi$ on $\R$, and the Fourier transform is exactly
\eqref{eq:original}. We use its Plancherel norm equality.

\begin{lemma}\label{lem:uncertainty}
For $h\in\Sch(\R)$ and integers $M,N\geq0$,
\begin{equation}\label{eq:uncertainty}
 \|h\|_2\leq2\left(4^M\|x^Mh\|_2
                         +4^N\|\xi^N\widehat h\|_2\right).
\end{equation}
\end{lemma}

\begin{proof}
Let $P$ and $Q$ be multiplication by the indicator of $[-1/4,1/4]$
on the physical and Fourier sides, respectively. The operator
$P\F^{-1}Q$ has kernel $1_{[-1/4,1/4]}(x)
e^{2\pi i x\xi}1_{[-1/4,1/4]}(\xi)$; its Hilbert--Schmidt norm is
$\sqrt{(1/2)(1/2)}=1/2$. Therefore, by Plancherel and the triangle
inequality,
\[
 \|h\|_2\leq\|Ph\|_2+\|(1-P)h\|_2
 \leq\tfrac12\|h\|_2+\|(1-Q)\widehat h\|_2+\|(1-P)h\|_2.
\]
On $|x|>1/4$, $1\leq4^M|x|^M$; the analogous statement holds in
frequency. Moving the half-norm to the left proves the assertion.
\end{proof}

Write the exact polynomial from \eqref{eq:R} as
$R_m(t)=\sum_{j=0}^m r_{m,j}t^j$. No coefficient is dropped in what
follows. For integers $c\geq3$ and $m\geq0$ put
\begin{align}
 E_{c,m}(F)&=\frac1{\sqrt2}\sum_{j=0}^m|r_{m,j}|
 \left(\frac{\zeta(c-1)}{\sqrt3}p_{c-1,j}(F)
                      +\zeta(c+1)p_{c+1,j}(F)\right),
                     \label{eq:E}\\
 E^J_{c,m}(F)&=\frac1{\sqrt2}\sum_{j=0}^m|r_{m,j}|
 \left(\frac{\zeta(c-1)}{\sqrt3}
                     \sum_{\ell=0}^j\binom j\ell p_{2-c,\ell}(F)
       +\zeta(c+1)\sum_{\ell=0}^j\binom j\ell p_{-c,\ell}(F)\right).
                     \label{eq:EJ}
\end{align}
For $a,b\in\N_0$ set, with $(b)_r=b!/(b-r)!$,
\begin{multline}\label{eq:A}
 A_{a,b}(F)=2\left[4^{b+3}E_{a+3,b}(F)\right.\\
 \left.{}+4^{a+3}(2\pi)^{b-a}
 \sum_{r=0}^{\min(a,b)}\binom ar(b)_rE^J_{b+3,a-r}(F)\right].
\end{multline}
Every $A_{a,b}$ is explicitly a finite nonnegative linear combination
of original $p_{b,j}$ seminorms.

\begin{proposition}\label{prop:inverse}
For $\phi\in V$, $F=\Theta\phi$, and $a,b\in\N_0$,
\begin{equation}\label{eq:L2-inverse}
 \|x^a\phi^{(b)}\|_2\leq A_{a,b}(F).
\end{equation}
In particular the ordinary Schwartz seminorms satisfy
\begin{equation}\label{eq:sup-inverse}
 \sup_{x\in\R}|x^a\phi^{(b)}(x)|
 \leq A_{a,b}(F)+aA_{a-1,b}(F)+A_{a,b+1}(F),
\end{equation}
where the middle term is omitted when $a=0$.
\end{proposition}

\begin{proof}
For any even Schwartz $\psi$, splitting the positive-axis integral at
$1$ gives
\[
 \|x^cD^j\psi\|_{L^2(\R)}
 \leq\sqrt2\left(\frac{p^+_{c-1,j}(\psi)}{\sqrt3}
                               +p^+_{c+1,j}(\psi)\right).
\]
Indeed, on $(0,1)$ the integrand's absolute value is at most
$xp^+_{c-1,j}$, whose squared integral is $(p^+_{c-1,j})^2/3$;
on $(1,\infty)$ it is at most $x^{-1}p^+_{c+1,j}$, whose squared
integral is $(p^+_{c+1,j})^2$. Evenness supplies the factor $\sqrt2$;
the square root of the sum is at most the sum of square roots.
Using the full polynomial $R_m$ and then
\eqref{eq:inverse-phi}--\eqref{eq:inverse-hat}, with $c-1\geq2$, proves
\begin{equation}\label{eq:U-bound}
 \|x^cR_m(D)\phi\|_2\leq E_{c,m}(F),\qquad
 \|x^cR_m(D)\widehat\phi\|_2\leq E^J_{c,m}(F).
\end{equation}

Apply Lemma~\ref{lem:uncertainty} to
$h=x^a\phi^{(b)}$, with $M=b+3$ and $N=a+3$. The first weighted
term is, by \eqref{eq:R},
\[
 x^Mh=x^{a+b+3}\phi^{(b)}=x^{a+3}R_b(D)\phi.
\]
The exact Fourier identities for the specified negative-sign kernel are
\[
 \widehat{\partial_x^b\phi}(\xi)=(2\pi i\xi)^b\widehat\phi(\xi),
 \qquad
 \widehat{x^a\psi}(\xi)=\left(-\frac1{2\pi i}\right)^a
                                      \partial_\xi^a\widehat\psi(\xi).
\]
Their Leibniz expansion therefore gives
\[
 \xi^{a+3}\widehat h(\xi)
 =\left(-\frac1{2\pi i}\right)^a(2\pi i)^b
 \sum_{r=0}^{\min(a,b)}\binom ar(b)_r
       \xi^{a+3+b-r}\widehat\phi^{(a-r)}(\xi).
\]
For each summand the final factor is exactly
$\xi^{b+3}R_{a-r}(D)\widehat\phi(\xi)$, where
$D=-\xi\partial_\xi$ on this axis. The scalar prefactor has modulus
$(2\pi)^{b-a}$. Substituting \eqref{eq:U-bound} into
\eqref{eq:uncertainty} proves \eqref{eq:L2-inverse} with exactly
\eqref{eq:A}.

For a Schwartz function $h$, integration from $-\infty$ to $x$ gives
$|h(x)|^2\leq2\|h\|_2\|h'\|_2$, whence
$\|h\|_\infty\leq\|h\|_2+\|h'\|_2$.
For $h=x^a\phi^{(b)}$ its derivative is
$a x^{a-1}\phi^{(b)}+x^a\phi^{(b+1)}$. The preceding $L^2$ estimate
and the triangle inequality yield \eqref{eq:sup-inverse}. The
seminorms on its left generate the full Schwartz topology; weights
$(1+|x|)^a$ are bounded by the finite sum
$\sum_{r=0}^a\binom ar|x|^r$ and conversely each $|x|^a$ is bounded by
$(1+|x|)^a$.
\end{proof}

\section{Closed image, the exact quotient, and the chain isomorphism}

\begin{theorem}\label{thm:closed}
The original map $\Theta:V\to\B$ is a continuous linear injection,
a topological isomorphism onto its range with the subspace topology,
and its range is closed. In particular
\begin{equation}\label{eq:closed}
 \overline{\Theta V}^{\,\B}=\Theta V.
\end{equation}
\end{theorem}

\begin{proof}
Continuity is \eqref{eq:theta-continuity}, injectivity is
\eqref{eq:mobius-inversion}, and inverse continuity is
\eqref{eq:sup-inverse}. Suppose $F_n=\Theta\phi_n$ converges in $\B$
to $F$. Applying \eqref{eq:sup-inverse} to
$\phi_n-\phi_m\in V$ shows that $(\phi_n)$ is Cauchy for every
Schwartz seminorm. Completeness and closedness of $V$ give
$\phi_n\to\phi\in V$. Continuity of $\Theta$ implies
$\Theta\phi_n\to\Theta\phi$ in $\B$, and its Hausdorff property gives
$F=\Theta\phi$. Thus the range is sequentially closed. Since $\B$ is
metrizable by its countable seminorm family, it is closed.
\end{proof}

The literature application reaches the same conclusion immediately
after Theorem~\ref{thm:topology}: Meyer's restricted closed embedding
is precisely $Z=\Theta/2$ with source $V=\HH_\cap$ and target
$\HH_-=\B$. Since $2V=V$, its image equals $\Theta V$ as a subset of
that identical target. The independent estimates above additionally
make the inverse continuity completely explicit in original seminorms.

Let $Q=\B/\Theta V$ with its quotient topology. It is Hausdorff:
if a coset is nonzero, its representative lies outside the closed
subspace $\Theta V$, and the quotient separates it from the zero
coset. It is also Fr\'echet. One direct completeness argument chooses
an increasing defining seminorm sequence for $\B$, uses the quotient
seminorms, selects from a quotient-Cauchy sequence a subsequence with
successive quotient distances below $2^{-n}$ in each of the first
$n$ seminorms, and lifts each successive difference to a vector whose
same finitely many seminorms are below $2^{1-n}$. The lifted series is
Cauchy in each seminorm, hence converges in $\B$ and projects to the
subsequence limit; the original Cauchy sequence has the same limit.
The quotient is metrizable by these countably many quotient seminorms.

The original algebraic exact sequence labelled (E5),
\[
 0\longrightarrow
 \overline{\Theta V}^{\,\B}/\Theta V\longrightarrow Q
 \longrightarrow\B/\overline{\Theta V}^{\,\B}\longrightarrow0,
\]
therefore has its first nontrivial term equal to zero. Its last map
is exactly
\begin{equation}\label{eq:E5}
 Q\xrightarrow{\ \cong\ }\B/\overline{\Theta V}^{\,\B},
 \qquad [F]_{\Theta V}\longmapsto[F]_{\overline{\Theta V}^{\,\B}},
\end{equation}
with inverse the same representative map. Both maps are continuous
because the two equivalence relations and quotient topologies are
identical. No new completion or representative change is involved.

To retain the operator factor at the chain level, define two-term
complexes in degrees $0,1$ by
\[
 C_\Theta=[V\xrightarrow{\Theta}\B],\qquad
 C_Z=[\HH_\cap\xrightarrow{Z}\HH_-].
\]
The explicitly typed chain isomorphism is
\begin{equation}\label{eq:chain}
 T^0=2\id:V\longrightarrow\HH_\cap,\qquad
 T^1=\id:\B\longrightarrow\HH_-.
\end{equation}
Indeed $ZT^0\phi=Z(2\phi)=\Theta\phi=T^1\Theta\phi$.
Its inverse has degree-zero map $\tfrac12\id$ and degree-one map
$\id$. Thus the induced degree-one quotient isomorphism is
$[F]\mapsto[F]$, and degree-zero cohomology is zero on both sides
by injectivity. These are the range complexes associated to the
operators; no different complex is attributed to Meyer.

For the original two-leg differential the equally literal diagram is
\[
 [V\oplus V\xrightarrow{\,\Theta-J\Theta\,}\B]
 \xrightarrow{\ (2\id\oplus2\id,\,\id)\ }
 [\HH_\cap\oplus\HH_\cap\xrightarrow{\,Z-JZ\,}\HH_-].
\]
The equality of differentials follows separately for each signed leg.
Its inverse scales both source coordinates by $1/2$ and fixes the
target. Since $J\Theta=\Theta\F$ on $V$, the source map
$(\phi_+,\phi_-)\mapsto\phi_+-\widehat\phi_-$ and the target identity
define a surjective chain map from the original two-leg complex to
$C_\Theta$. Its kernel is the degree-zero graph
$\{(\widehat\psi,\psi):\psi\in V\}$ with zero degree-one term:
both inclusions follow directly from
$\phi_+-\widehat\phi_-=0$. In particular the two-leg degree-one image
is also precisely $\Theta V$, not a different closure.

Finally $D(V)\subset V$: it preserves evenness, $D\phi(0)=0$, and
integration by parts gives $\int_\R D\phi=\int_\R\phi=0$.
It acts continuously on $\B$ because $p_{b,j}(DF)=p_{b,j+1}(F)$,
and termwise differentiation gives $D\Theta=\Theta D$. Thus the
one-leg chain isomorphism commutes with $D$ in both degrees.
The signs for the two-leg complex are also explicit. Integration by
parts and differentiation of the Fourier kernel give
\[
 \F D\phi=(1+\xi\partial_\xi)\widehat\phi
           =(1-D)\F\phi,
 \qquad \F(1-D)\phi=D\F\phi.
\]
Consequently its source generator is $D\oplus(1-D)$ and its target
generator is $D$: by \eqref{eq:J},
\[
 D(\Theta\phi_+-J\Theta\phi_-)
    =\Theta D\phi_+-J\Theta(1-D)\phi_-.
\]
The two-leg scaling isomorphism commutes with these generators.
So does the projection to the one-leg complex, because
$\phi_+-\widehat\phi_-$ is sent under the source generator to
$D\phi_+-\F(1-D)\phi_-=D(\phi_+-\widehat\phi_-)$.
No common diagonal generator on the two source legs and no sign
change to $x\partial_x$ is silently substituted.

\section{The exact scope for subspaces, finite cutoffs, and Hilbert norms}

\begin{proposition}\label{prop:subspace}
For every linear subspace $W\subseteq V$, with closure taken in the
indicated original spaces,
\begin{equation}\label{eq:W}
 \overline{\Theta W}^{\,\B}
       =\Theta\left(\overline W^{\,V}\right).
\end{equation}
The restriction of $\Theta$ induces a topological linear isomorphism
\[
 \overline W^{\,V}/W
       \longrightarrow\overline{\Theta W}^{\,\B}/\Theta W,
 \qquad[\psi]\longmapsto[\Theta\psi],
\]
with the respective quotient topologies, whether or not the quotients
are Hausdorff. In particular every finite-dimensional $W$ has closed
theta image in $\B$.
\end{proposition}

\begin{proof}
The range $R=\Theta V$ is closed in $\B$ and $\Theta:V\to R$ is a
homeomorphism. Therefore the closure of $\Theta W$ in $\B$ lies in
$R$ and equals its closure in $R$. A homeomorphism carries closures
to closures, proving \eqref{eq:W}. Its restrictions to the closed
subspaces in that equality are inverse homeomorphisms. They carry
the equivalence classes modulo $W$ and $\Theta W$ to each other,
so the quotient universal property proves both induced maps
continuous and inverse.

For the last assertion, finite-dimensional subspaces of the present
Hausdorff locally convex space are closed. Here is a direct proof.
For a basis $w_1,\ldots,w_d$ of $W$, the map
$T:\C^d\to V$, $c\mapsto\sum c_iw_i$, is continuous and injective.
For each vector on the Euclidean unit sphere choose a continuous
seminorm positive on its image. Compactness gives finitely many
such seminorms whose sum $p$ is positive on every image on the
sphere; continuity and compactness give $p(Tc)\geq\delta|c|$ for
some $\delta>0$. Thus a convergent sequence $Tc_n$ in $V$ has
Cauchy coordinates, with limit $c$, and its limit is $Tc\in W$.
Metrizability gives closedness of $W$. Now apply \eqref{eq:W}.
\end{proof}

Consequently for any programme subspace denoted $W_h$, formula
\eqref{eq:W}, not an unproved replacement of $W_h$ by $V$, is the
exact available assertion. It computes its closure defect as the
source defect $\overline{W_h}^{\,V}/W_h$. For each fixed finite
polynomial source span $W_L$, its image is closed in $\B$ by the
finite-dimensional assertion. For a union of these spans the precise
formula remains
\[
 \overline{\bigcup_L\Theta W_L}^{\,\B}
    =\Theta\left(\overline{\bigcup_L W_L}^{\,V}\right)
\]
when the spans are increasing, so their union is a linear subspace.
No assertion that this source closure is all of $V$ is made.

The inclusion into the original Hilbert space
$\HH=L^2(\R_{>0},dx)$ is a different, explicitly continuous map:
\begin{equation}\label{eq:Hilbert}
 \|F\|_{\HH}^2
 \leq p_{0,0}(F)^2\int_0^1dx
          +p_{1,0}(F)^2\int_1^\infty x^{-2}\,dx
 =p_{0,0}(F)^2+p_{1,0}(F)^2.
\end{equation}
It is injective since a continuous function zero almost everywhere
is zero everywhere. Its inverse on its image is not continuous.
Indeed fix $0\ne\eta\in C_c^\infty((1,2))$ and put
$F_n(x)=n^{-1}\eta(x-n)$. Each $F_n$ belongs to $\B$;
$\|F_n\|_{\HH}=n^{-1}\|\eta\|_2\to0$, whereas for a fixed
$t_0\in(1,2)$ with $\eta(t_0)\ne0$,
\[
 p_{2,0}(F_n)\geq n^{-1}(n+t_0)^2|\eta(t_0)|\longrightarrow\infty.
\]
Thus Hilbert convergence does not imply convergence in the original
target topology. In particular the closed-image theorem does not
assert Hilbert closedness, alter any finite Hilbert Gram matrix, or
replace a calculation of Hilbert closure. It removes exactly the
closure defect in (E5) for the original full source $V$ and the
original target seminorm topology.

\end{document}
